\documentclass{article}

\usepackage{amsmath,amssymb}
\usepackage{xurl}
\usepackage[hidelinks]{hyperref}
\setlength{\emergencystretch}{2em}

% Shared commands for notes in docs/paper-gaps/.

\DeclareUnicodeCharacter{2080}{\ifmmode{}_{0}\else\textsubscript{0}\fi}
\DeclareUnicodeCharacter{2081}{\ifmmode{}_{1}\else\textsubscript{1}\fi}
\DeclareUnicodeCharacter{2082}{\ifmmode{}_{2}\else\textsubscript{2}\fi}
\DeclareUnicodeCharacter{2099}{\ifmmode{}_{n}\else\textsubscript{n}\fi}

\newcommand{\Lean}{\textsc{Lean}}
\ExplSyntaxOn
\NewDocumentCommand{\leanid}{m}{
  \ifmmode
    \textsf{\detokenize{#1}}
  \else
    \group_begin:
    \urlstyle{sf}
    \tl_set:Nn \l_tmpa_tl {#1}
    \tl_replace_all:Nnn \l_tmpa_tl {\_} {_}
    \exp_args:NV \path \l_tmpa_tl
    \group_end:
  \fi
}
\ExplSyntaxOff
\providecommand{\tr}{\operatorname{tr}}
\newcommand{\tnleanIssueBase}{https://github.com/LionSR/TNLean/issues/}
\newcommand{\tnleanPrBase}{https://github.com/LionSR/TNLean/pull/}
\newcommand{\tnleanDocBase}{https://sirui-lu.com/TNLean/docs/}
\newcommand{\ghissue}[1]{\href{\tnleanIssueBase#1}{GitHub issue \##1}}
\newcommand{\ghpr}[1]{\href{\tnleanPrBase#1}{PR \##1}}
\newcommand{\doclink}[1]{\href{\tnleanDocBase#1.html}{\path{#1}}}

% Dirac notation and the identity operator, as in blueprint/src/macros/common.tex.
% \providecommand keeps note-local definitions authoritative.
\usepackage{dsfont}
\providecommand{\ket}[1]{|#1\rangle}
\providecommand{\bra}[1]{\langle #1|}
\providecommand{\braket}[2]{\langle #1 | #2 \rangle}
\providecommand{\ketbra}[2]{|#1\rangle\!\langle #2|}
\providecommand{\Id}{\mathds{1}}
\providecommand{\one}{\mathds{1}}
\providecommand{\C}{\mathbb{C}}
\providecommand{\MN}[1]{M_{#1}(\C)}
\providecommand{\MD}{M_D(\C)}

% External referencing for paper sources.  The build helper writes sanitized
% auxiliary files under build/paper-aux/ and excludes duplicated source labels.
\usepackage{xr}

% Import a generated paper auxiliary file with a caller-supplied label prefix.
% The candidate paths support builds from docs/paper-gaps/, the repository root,
% and build/paper-aux/ itself.
\newcommand{\importpaperaux}[2]{%
  \IfFileExists{../../build/paper-aux/#2.aux}{%
    \externaldocument[#1]{../../build/paper-aux/#2}%
  }{%
    \IfFileExists{build/paper-aux/#2.aux}{%
      \externaldocument[#1]{build/paper-aux/#2}%
    }{%
      \IfFileExists{#2.aux}{%
        \externaldocument[#1]{#2}%
      }{}%
    }%
  }%
}

% General external reference macros
\newcommand{\paperref}[2]{\ref{#1-#2}}
\newcommand{\papereqref}[2]{(\ref{#1-#2})}

% CPSV16 source (1606.00608 / MPDO-22-12-17-2.tex) external document and shortcuts
\importpaperaux{cpsv16-}{1606.00608/MPDO-22-12-17-2-tnlean}
\newcommand{\cpsvref}[1]{\paperref{cpsv16}{#1}}
\newcommand{\cpsveqref}[1]{\papereqref{cpsv16}{#1}}

% Verdict marker: \gapnote{<kind>}{<status>}. Kind is the policy
% classification (clarification, local-correction, scope-restriction,
% unfaithful, false-source, open-gap); status is open, wip (elimination
% underway), resolved, or historical. Machine-read by the site index;
% prints a verdict line.
\newcommand{\gapnote}[2]{\par\noindent\textbf{Verdict:} #1 (#2).\par}


\title{SIC--POVM Equality Families: the One-Dimensional Degeneracy}
\author{}
\date{2026-08-21}

\begin{document}
\maketitle
\gapnote{local-correction}{resolved}

\section{Source passage}

Let $P_1,\ldots,P_n\in\MN{d}$ be positive semidefinite matrices satisfying
$\tr(P_i^2)=1$, where $d\le n$. Chapter~2 of
\cite{Wolf2012Quantum}, Proposition ``SIC POVMs'', proves the bound
\begin{align}
  \sum_{i\ne j}|\tr(P_iP_j)|^2
  \ge \frac{n(n-d)^2}{(n-1)d^2}.
  \label{eq:wolf-sic-bound}
\end{align}
At equality, the argument gives
\begin{align}
  P_i^2&=P_i, & \operatorname{rank}(P_i)&=1, \\
  \sum_iP_i&=\frac nd\Id, &
  \tr(P_iP_j)&=\frac{n-d}{(n-1)d}\qquad(i\ne j).
  \label{eq:wolf-sic-equality}
\end{align}
The source then asserts, at lines 816--823 of the local Chapter~2 source file,
that the matrices in an equality family are linearly independent.

\section{The deviation}

The concluding assertion fails when $d=1$ and $n>1$. In one dimension the
unit-purity and positivity hypotheses force every $P_i$ to be the matrix
$[1]$. This family attains equality in~\eqref{eq:wolf-sic-bound}: both sides
are $n(n-1)$. It is nevertheless linearly dependent as soon as it contains
more than one member.

The coefficient calculation in the source also isolates the precise
degeneracy. If $\sum_i c_iP_i=0$, taking the trace and then pairing with
$P_j$ gives
\begin{align}
  \sum_i c_i&=0, \\
  \left(1-\frac{n-d}{(n-1)d}\right)c_j&=0.
\end{align}
For $n>1$, the prefactor vanishes exactly when $d=1$. Thus an equality family
is linearly independent under the corrected condition $d>1$. When $n=1$,
Equation~\eqref{eq:wolf-sic-bound} is not stated because of its denominator,
but the singleton family is linearly independent whenever its sole matrix has
unit purity.

\section{Formalized interpretation}

The formalization separates three statements. The inequality
\leanid{Matrix.sicPOVM_offDiag_overlap_sq_bound} and its equality criterion
\leanid{Matrix.sicPOVM_offDiag_overlap_sq_eq_iff} retain the source's
nondegenerate hypotheses $n\ge2$ and $1\le d\le n$; in particular, their
validity at $d=1$ is preserved. The corrected consequence
\leanid{Matrix.sicPOVM_linearIndependent_of_overlap_bound_eq} assumes $d\ge2$.
The complementary singleton statement is
\leanid{Matrix.singleton_linearIndependent_of_trace_sq_eq_one}.

Consequently, the combined nondegeneracy condition for linear independence is
$d>1$ or $n=1$, while the equality theorem itself is not artificially
restricted.

\section{Verdict}

\textbf{Verdict: resolved local fix (one-dimensional equality family).}
Wolf's overlap bound and equality characterization are correct in one
dimension. Only the final linear-independence consequence requires the
qualification $d>1$, with the singleton case treated separately.

\bibliographystyle{alpha}
\bibliography{references}

\end{document}
